\scanpage{304}
\mspage{\textcircled{3}}
\[
 K^{[m]}{}_{[n]}=K(\Delta^m,\Delta^n)\in A
\]
De $=$, on note
\[
 \begin{aligned}
 K_*{}^*&\quad\text{les objets de }(A^*)_*\simeq\underline{\mathrm{Hom}}(\mathbb{\Delta}^o\times\mathbb{\Delta},A)\\
 K^{**}&\quad\text{les objets de }(A^*)^*\simeq\underline{\mathrm{Hom}}(\mathbb{\Delta}^{\unsure{o}}\times\mathbb{\Delta},A)\\
 K_{**}&\quad\text{les objets de }A_{**}=\underline{\mathrm{Hom}}(\mathbb{\Delta}^o\times\mathbb{\Delta}^o,A)
 \end{aligned}
\]

On fera attention que si on compose avec \inserted{l'un} des foncteurs de symétrie \deleted{$\unsure{(m,n)\mapsto(n,m)}$} \inserted{$(\Delta^m,\Delta^n)\longmapsto(\Delta^n,\Delta^m)$} dans $\mathbb{\Delta}\times\mathbb{\Delta}$ ou dans $\mathbb{\Delta}^o\times\mathbb{\Delta}^o$, ou entre $\mathbb{\Delta}\times\mathbb{\Delta}^o$ et $\mathbb{\Delta}^o\times\mathbb{\Delta}$), on n'utilise pas la $=$ lettre pour désigner un \deleted{complexe} \inserted{objet} (tel que $K^*{}_*$ ou $K^{**}$ disons) ou son transformé (qui devrait alors s'écrire $K_*{}^*$, ou encore une fois $K^{**}$), car cela conflicterait avec la convention plus haut de la ``montée et descente des indices'' $K_*{}^*\simeq K^*{}_*{}^\vee$ dans le premier cas, et \deleted{\unsure{serait}} amènerait à écrire $K^{[m][n]}=K^{[n][m]}$ dans le \unsure{deuxième} ! On utilisera au besoin un symbole ``de symétrie'' tel que $\sigma$, pour dénoter le symétrique de $K$ ($=K^*{}_*,K^{**}$) par $ {}^{\sigma}K$ ($=({}^{\sigma}K)_*{}^*,({}^{\sigma}K)^{**}$) \deleted{\unsure{de|dans}}\ldots\ D'ailleurs, \deleted{Ces considérations et notations} si on a un objet $K^*{}_*\inserted{{}\in\mathrm{Ob}\,(A_{\unsure{*}}^{*})^{*}_{\deleted{\unsure{*}}}}$ disons, alors pour \unsure{tt} $n\geq0$, on peut considérer
\[
 \deleted{\unsure{K\longmapsto}}\quad K^*{}_{[n]}\overset{\text{déf}}{=}(m\longmapsto K^{[m]}{}_{[n]})
\]
comme un objet de $A^*$ ; le foncteur $[n]\longmapsto K^*{}_{[n]}$ de $A^*{}_*$ n'est autre que $ {}^{\sigma}K$.

On étend ces considérations et notations au cas de symboles $K$ à plus de deux indices, pour désigner des objets de catégories du
% archive sheet number: 305
